\scanpage{197}
\mspage{71}
\[
 \deleted{i_\alpha\in\bigcap_{\beta\leq\alpha}I_\beta\quad\unsure{\text{pour}}\ .}
\]
\deleted{On choisit $i_0$ ($0$ élément initial de $A$) dans $I_0$ quelconque, et on}

que pour tout $\alpha\in A$, l'ensemble
\[
 J_\alpha=\bigcap_{\beta<\alpha}I_\beta
\]
est cofinal. Si on sait que $J_\alpha$ est cofinal, il en est de $=$ de $J_{\alpha+1}=J_\alpha\cap I_\alpha$, par 5.5. Donc on est ramené à prouver que $J_\alpha$ est cofinal si $\alpha$ est un ordinal limite, donc à prouver le lemme 5.15. pour une suite \inserted{transfinie} \uline{décroissante} de parties \deleted{\unsure{dans un}} $\unsure{\pi|R|K}$-cofinales, \inserted{$J_\beta$, $(\beta<\alpha)$ \unsure{quand} $\operatorname{card}\alpha\leq\pi$.} \deleted{\unsure{de}} \inserted{telle que $J_\beta=\bigcap_{\beta'<\beta}J_{\beta'}$, \deleted{\unsure{pour}} \unsure{quand} $\beta$ est limite.} \inserted{\unsure{A} prouver} que l'intersection est $\pi$-cofinale --- et on est réduit à prouver qu'elle est non vide. On va construire \unsure{par} réc. transfinie une suite \uline{croissante} d'éléments
\[
 i_\beta\in J_\beta\qquad(\beta\leq\alpha)
\]
Si $i_\beta$ est construit, comme $J_{\beta+1}$ est cofinal, il existe $i_{\beta+1}\geq i_\beta$ dans $J_{\beta+1}$, OK. Si $\beta$ \unsure{est} un ordinal limite, on prend
\[
 i_\beta=\mathop{\mathrm{b.\!sup}}\limits_{\beta'<\beta} i_{\beta'}, 
\]
$\unsure{i'_\beta|i_{\beta'}}$ est dans \deleted{\unsure{les}} chaque $I_{\beta_0}$ pour $\beta_0<\beta$,
% archive sheet number: 196
